% Supplemental F.1 audit. All four supplied proofs were read before this check.
\subsection{The remaining background checks: R7--R10}

\begin{proposition}[R7, scans 9:16--18]\label{prop:DE-R7}
Let $A_0/k$ be an abelian variety in characteristic $p$.  A base-change-compatible, functorial ordinary $k$-stratification on $H^1_{\dR}(A/\Lambda)$ for every infinitesimal deformation of $A_0$, horizontal for the relative operators
$F:H^{(p)}\to H$ and $V:H\to H^{(p)}$, exists if and only if $A_0$ is ordinary.  In that case it is unique and preserves the canonical splitting
\[
 H=\omega_{A/\Lambda}\oplus\im F.
\]
Thus the Hodge filtration is horizontal, as asserted in the manuscript.
\end{proposition}
\begin{proof}
The formally smooth universal unpolarized deformation space is $\operatorname{Spf}k[[t_1,\ldots,t_{g^2}]]$; this follows from the Barsotti--Tate deformation theorem and Serre--Tate, in \cite[Corollary~4.8(i)]{Illusie1985BT} and \cite[Theorem~1.2.1]{Katz1981ST}.  Pulling a formal stratification back by $(0,t)$ trivializes its underlying module over this complete local ring.  Horizontality then identifies every iterate of $F$ and $V$ with the corresponding operator on the closed fibre.  Their geometric stable ranks cannot change.

For clarity, the ordinary-deformation input can be verified from an opened explicit display, without an unlocated assertion about density.  After extending the residue field to an algebraic closure, split off the \etale part of $A_0[p^\infty]$.  For the remaining formal group of height $h$ and dimension $d$, write Oort's covariant display matrix as
\[
 \begin{pmatrix}A&B\\C&D\end{pmatrix}\in\GL_h(W(k)),\qquad
 F=\begin{pmatrix}A&pB\\C&pD\end{pmatrix}\sigma.
\]
The construction in \cite[\S\S1.9, 1.12--1.14]{Oort2000} replaces the first matrix by
$\left(\begin{smallmatrix}A+TC&B+TD\\C&D\end{smallmatrix}\right)$, for independent formal parameters $T$ of size $d\times(h-d)$.  Modulo $p$, the column matrix $\binom A C$ has rank $d$.  Choose $T$ so that $A+TC$ is invertible: $A$ is injective on $\ker C$, and one prescribes $T$ on $\im C$ so that the images of a complementary subspace fill a complement to $A(\ker C)$.  Hence $\det(A+TC)$ is a nonzero polynomial and is nonzero at the generic point of the formal parameter space.  In this \emph{covariant} convention, the maximal stable rank of $F$ modulo $p$ is the dimension $d$ of the group.  That maximum is attained here, equivalently the group is ordinary.  Adding back the fixed \etale part and applying Serre--Tate gives an ordinary generic deformation of $A_0$.  If $A_0$ were nonordinary this would contradict the constant stable ranks forced by the stratification.  All arguments may be checked after the faithfully flat extension of the residue field.

For an ordinary deformation, $F$ kills $\omega^{(p)}$, and its induced map on $(H/\omega)^{(p)}$ is invertible.  Thus $\im F$ is a complement to $\omega$.  The two summands are the vector bundles associated with the appropriate \etale $\Fp$-local systems of $A[p]$ and $A^*[p]$, with their invertible $F$- or $V$-structures.  Nilpotent invariance of the \etale site supplies their canonical ordinary stratifications.  They are functorial and horizontal for both operators.

Any other horizontal stratification preserves $\im F$ and also $\im V=\omega^{(p)}$.  Relative Frobenius on the formally smooth universal base, and on its double product, is faithfully flat; therefore preservation of $\omega^{(p)}$ implies preservation of $\omega$.  On either summand, an automorphism comparing two such stratifications is a matrix $G$ with $G\equiv1$ modulo the diagonal ideal $J$ and $G=G^{(p)}$ in a basis of its \etale $\Fp$-structure.  Iteration gives $G\equiv1\pmod{J^{[p^n]}}$ for every $n$, hence $G=1$.  Pullback from the universal base proves uniqueness for every infinitesimal deformation.
\end{proof}

\begin{proposition}[R8, scans 9:5--8]\label{prop:DE-R8}
Let $k$ be perfect, $W=W(k)$, $W_m=W/p^m$, and let
\[
 \mathcal D=W[F,V],\qquad F\lambda=\sigma(\lambda)F,\quad
 \lambda V=V\sigma(\lambda),\quad FV=VF=p.
\]
For $\mathcal D_n=\mathcal D/(F^{n+1},V^{n+1})$, the constant-term pairing gives an isomorphism of $\mathcal D$-bimodules
\[
 \mathcal D_n\xrightarrow{\sim}
 \Hom_{W_{n+1}}({}_{W_{n+1}}\mathcal D_n,W_{n+1}),\qquad
 y\longmapsto[x\longmapsto\phi(xy)].
\]
The coefficient ring $W_n$ on the handwritten page must be replaced by $W_{n+1}$.  With this correction, the two dualizing functors on modules killed by powers of both $F$ and $V$ agree.
\end{proposition}
\begin{proof}
The normal form is
\[
 \mathcal D_n=W_{n+1}\oplus
 \bigoplus_{i=1}^nW_{n+1-i}F^i\oplus
 \bigoplus_{i=1}^nV^iW_{n+1-i}.
\]
In particular its $W$-length is $(n+1)^2$ and $p^n1\ne0$, which already rules out the proposed $W_n$-linear formula.  The constant-term map is left and right $W$-linear and satisfies
\[
 \phi(Fz)=\sigma\phi(zF),\qquad
 \phi(zV)=\sigma\phi(Vz).
\]
The constant summand pairs with itself by multiplication.  The $F^i$ and $V^i$ summands pair only with each other, through multiplication by $p^i$ and a bijective $\sigma^i$ twist.  They pair perfectly because
$\Hom_{W_{n+1}}(W_{n+1-i},W_{n+1})=p^iW_{n+1}$.

Here are the actions, to remove the two transcription slips and the possible ambiguity of the word ``bimodule.''  For a left $W$-linear functional $u$, set
\[
 (a u)(x)=u(xa),\quad
 (u\lambda)(x)=u(\lambda x),\quad
 (uF)(x)=\sigma^{-1}u(Fx),\quad
 (uV)(x)=\sigma u(Vx).
\]
They satisfy the Dieudonn\'e relations and make the pairing a bimodule map.  Thus the right $V$ action uses $V$, not $F$, and the left $V$ action is precomposition with right multiplication by $V$.

The injections $\mathcal D_n\to\mathcal D_{n+1}$ given by multiplication by $p$ are well defined, since $pF^{n+1}=VF^{n+2}$ and $pV^{n+1}=FV^{n+2}$ in the latter quotient.  They identify the annihilator of $F^r,V^r$ in $I=\varinjlim_n\mathcal D_n$ with $\mathcal D_{r-1}$.  Therefore, if $F^rM=V^rM=0$, adjunction and the pairing identify
\[
 \Hom_{\mathcal D}(M,I)
 =\Hom_{\mathcal D}(M,\mathcal D_{r-1})
 \simeq\Hom_{W_r}(M,W_r).
\]
The anti-involution fixing $W$ and exchanging $F,V$ converts the right module into the usual left Dieudonn\'e dual.  This independently verifies the supplied argument, including its indices.
\end{proof}

\begin{proposition}[R9, scans 6:22--27]\label{prop:DE-R9}
No ordinary stratification of $H^1_{\dR}$ can be functorial for all elliptic schemes in characteristic $p$ and commute with all base changes.  The three corresponding assertions on independence of a lift, recovery of the classical Dieudonn\'e module, and classification of lifts by their Hodge filtration hold on nilpotent divided-power thickenings, with $p$ locally nilpotent.
\end{proposition}
\begin{proof}
For an explicit obstruction take $k=\overline{\mathbb F}_3$ and the Legendre family
\[
 y^2=x(x-1)(x-\lambda),\qquad \lambda=2+t,
\]
on the formal neighbourhood of $t=0$.  The Hasse polynomial is $1+\lambda=t$.  More generally, for every odd prime, put $m=(p-1)/2$.  The Legendre Hasse polynomial, up to a nonzero scalar depending on the differential convention, is
\[
 H_p(\lambda)=\sum_{i=0}^m\binom mi^2\lambda^i.
\]
It has positive degree, $H_p(0)=1$, and
$H_p(1)=\binom{p-1}{m}\ne0\pmod p$.  Over an algebraic closure choose a root $\lambda_0$, necessarily different from $0,1$, and set $\lambda=\lambda_0+t$.  The discriminant $16\lambda^2(1-\lambda)^2$ is a unit at $t=0$, while the generic Hasse invariant is nonzero.  For $p=2$, use instead
\[
 y^2+txy+y=x^3;
 \qquad\Delta=t^3+1,\qquad\operatorname{Hasse}=t.
\]
These give a smooth family in every characteristic with supersingular closed fibre and ordinary generic fibre.  Thus $F^2=0$ on the closed-fibre de Rham cohomology and $F^2\ne0$ generically.  Functoriality for relative Frobenius would make $F$ horizontal.  Pullback of an ordinary formal stratification by $(0,t)$ would then propagate $F^2=0$ to the whole formal neighbourhood, a contradiction.  This is the argument of the manuscript, also published in \cite[\S3.5]{Grothendieck1968Crystals}, with fully specified families.  The retained program \texttt{code/R7\_R10\_checks.py} verifies the two displayed discriminants and the Legendre coefficient formula by exact symbolic arithmetic; the general binomial identity above supplies the proof for all odd primes.

For the positive assertion write $D(G)$ for the \emph{contravariant} Dieudonn\'e crystal, so that
\[
 0\longrightarrow\omega_G\longrightarrow D(G)_S
       \longrightarrow\operatorname{Lie}(G^*)\longrightarrow0.
\]
For $G=A[p^\infty]$, its evaluation on a PD thickening is canonically $H^1_{\dR}$ of any abelian lift, with the displayed Hodge filtration.  The published locators are \cite[\S2.2, Theorem~2; \S3, Proposition~2 and (3.1)]{BerthelotMessing1979}.  In the original universal-extension construction, see \cite[Chapter~I, Proposition~4.1.7; Chapter~II, (9.2), (13.5), Theorem~15.3]{MazurMessing1974}; the last theorem identifies evaluation over a perfect field with their classical contravariant module.  If instead one uses $M_k(G)=\Hom(G,CW)$ with the precise convention of \cite[\S4.3, Theorem~4 and Remark~1]{BerthelotMessing1979}, that comparison is $\sigma$-semilinear: its $W(k)$-linear form is $\sigma^*M_k(G)\simeq\Gamma D(G)$.  This harmless Frobenius twist must be recorded rather than silently identifying the two conventions.

The lifting equivalence itself is \cite[Chapter~V, Theorem~1.6 for Barsotti--Tate groups; Theorem~1.10 for abelian schemes]{Messing1972}: admissible filtrations are locally direct summands with locally free quotient lifting the special-fibre filtration.  In Messing's original covariant convention use $\operatorname{Lie}E(G)$, or apply it to $G^*$ to obtain the contravariant convention above.  The local-liftability qualification for Barsotti--Tate groups in that book is supplied by \cite[Th\'eor\`eme~4.4 and Corollary~4.7]{Illusie1985BT}; for abelian schemes it is already \cite[Chapter~IV, Lemma~2.8.1]{Messing1972}.

Here ``nilpotent divided powers'' is a condition on the divided-power filtration, not merely nilpotence of the underlying ideal.  The distinction explains why the negative ordinary-stratification argument does not contradict the positive crystalline statement: in characteristic $p$, a PD element $t$ satisfies $t^p=p!\gamma_p(t)=0$, whereas the ordinary formal neighbourhood permits arbitrary powers of $t$.
\end{proof}

\begin{proposition}[R10, scans 7:16--17]\label{prop:DE-R10}
If $X/k$ is smooth and proper and $k$ is algebraically closed of characteristic $p$, then
\[
 H^i_{\mathrm{str}}(X/k,\OO)
 \simeq H^i(X,\OO_X)_{\ss}
 \simeq H^i_{\et}(X,\Fp)\otimes_{\Fp}k.
\]
For $W=W(k)$, the no-PD infinitesimal cohomology $H^i((X/W)_{\inf},\OO)$ vanishes for $i>\dim X$.  Thus the ordinary abelian example in R10 is valid; the vanishing does not require ordinarity.
\end{proposition}
\begin{proof}
For smooth $X$, formal local lifting identifies ordinary stratifications and infinitesimal crystals.  Frobenius-kernel diagonals are cofinal in the ordinary infinitesimal diagonals, and Frobenius on $X$ is faithfully flat.  Their descent description computes the stable part of proper coherent cohomology.  On a finite-dimensional $k$-space with semilinear Frobenius, the stable image is the scalar extension of its $F=1$ subspace; $F-1$ is surjective.  Artin--Schreier consequently identifies this space with $H^i_{\et}(X,\Fp)\otimes k$.  This is also the $W_1$ case of \cite[Corollary~4.7]{Ogus1975} and of Theorem~\ref{thm:DE-inf}, whose compatible local-embedding argument removes the global-embedding hypothesis.

For $i>\dim X$, coherent cohomology is zero, so the mod-$p$ \etale group is zero.  Induction on $n$ using the coefficient exact sequences gives $H^i_{\et}(X,\mathbb Z/p^n)=0$.  Passing to the inverse limit gives the assertion by Theorem~\ref{thm:DE-inf}; its Mittag--Leffler verification handles the preceding cohomological degree.  This is integral, including torsion, rather than a rational slope-zero assertion alone.
\end{proof}
